% !TEX root = ../../main.tex
\begingroup
\newcommand{\notationmeaning}[1]{{\raggedright#1\par}}
\chapter{Notation}
\label{ch:appendix:notation}
\raggedbottom

This appendix explains the symbols and reading conventions used in the book.
Meanings are local to the stated space, map, or matrix.
\autoref{ch:appendix:formula-sheet} collects the corresponding formulas.

\section{Scalars, spaces, and sets}
\label{sec:appendix:notation-spaces}

\begin{table}[htbp]
  \centering
  \renewcommand{\arraystretch}{1.2}
\begin{tabular}{@{}p{0.23\linewidth}p{0.72\linewidth}@{}}
  \toprule
  Symbol & Meaning \\
  \midrule
  \(\mathbb F\) & \notationmeaning{The field of scalars.} \\
  \(\R,\C\) & \notationmeaning{The real and complex numbers, respectively.} \\
  \(\mathbb F^n\) & \notationmeaning{Coordinate vectors with \(n\) entries
    in \(\mathbb F\).} \\
  \(V,W\) & \notationmeaning{Vector spaces; \(v,w\) usually denote
    their vectors.} \\
  \(V^n\) & \notationmeaning{Ordered \(n\)-tuples of vectors in \(V\),
    the inputs to an \(n\)-input multilinear function.} \\
  \(\dim V\) & \notationmeaning{The dimension of \(V\).} \\
  \(\operatorname{span}(S)\) & \notationmeaning{All finite linear
    combinations of vectors in \(S\).} \\
  \bottomrule
\end{tabular}
  \caption{Scalars, vector spaces, and span.}
  \label{tab:appendix:notation-spaces}
\end{table}

\autoref{tab:appendix:notation-spaces} uses the locally specified scalar
field. \autoref{ch:linear-algebra-i:determinants}
uses real vector spaces. An ordered basis written with an ordinary letter
\(F\) is distinct from the scalar field \(\mathbb F\).

In set notation, \(v\in V\) means that \(v\) is an element of \(V\),
while \(S\subseteq V\) means that \(S\) is a subset of \(V\).
The symbols \(\varnothing\) and \(\{0\}\) denote the empty set and
the set containing the zero vector, respectively. In a set such as
\(\{v\in V\mid L(v)=0\}\), the vertical bar means ``such that'';
a colon is also used for this purpose.

The symbol \(0\) can denote a scalar, vector, or zero matrix. Subscripts
in \(0_V\) and \(0_W\) specify the space containing the zero vector.
Thus \(\ker L=\{0_V\}\) describes a subspace with one element,
whereas \(\dim\ker L=0\) describes its dimension.

\section{Coordinates and matrix indices}
\label{sec:appendix:notation-indices}

% The supplied September 24 handout, pp. 18--19, and Homework 5 use
% upper row indices and lower column indices, consistently with the book.
For a matrix \(M=(M^i_j)\), the index positions have the following meaning:
\[
  M^i_j=\text{entry in row \(i\), column \(j\)}.
\]
The \emph{upper} index selects the row; the \emph{lower} index selects
the column. In notation that places both indices below the letter, the
same entry would be written \(m_{ij}\), with the row index first.
For example, \(M^2_3\) is the entry in the second row and third column.
Here the superscript is an index, not a power.

The notation \(M_j\) denotes the entire \(j\)-th column. Its entries are
\[
  M_j=\begin{bmatrix}M^1_j&\cdots&M^r_j\end{bmatrix}^{T}
  \qquad\text{for a matrix with \(r\) rows.}
\]
The superscript \(T\) means transpose, so the displayed row becomes a
column. Similarly, \(v=(v^1,\ldots,v^n)^T\) is a coordinate column;
\(v^i\) is its \(i\)-th scalar coordinate. By contrast, \(v_j\) usually
labels the \(j\)-th vector in a list. In the expression
\[
  v_j=\sum_{i=1}^n b^i_j e_i,
\]
the lower index \(j\) selects the vector, and the upper index \(i\)
selects its coordinate in the basis direction \(e_i\).

The size \(r\times c\) always means \(r\) rows and \(c\) columns.
Both \(\mathrm M_{r\times c}(\mathbb F)\) and
\(\mathbb F^{r\times c}\) denote matrices of this size with scalar
entries. The notation \(\mathrm M_{r\times c}(V)\) instead allows
entries in a vector space \(V\). Read dimension labels \(m,n\) from the
stated matrix size or map domain and codomain.

\subsection{Summation conventions}

An explicit sum such as \(\sum_{i=1}^n f_iM^i_j\) leaves \(j\) fixed
and sums over \(i\). Where Einstein notation is used, an index repeated
once above and once below is summed over its stated range. Thus
\[
  f_iM^i_j=\sum_{i=1}^n f_iM^i_j,
  \qquad
  e_ja^j=\sum_{j=1}^m e_ja^j.
\]
An index that is not summed remains free. For instance, \(j\) in
\(L(e_j)=f_iM^i_j\) identifies the basis vector whose image is being
computed. These conventions are used in
\autoref{sec:linear-algebra-i:change-of-basis-for-functions-and-maps}.

\section{Ordered bases and maps}
\label{sec:appendix:notation-bases-and-maps}

An ordered basis is written as \(E=(e_1,\ldots,e_m)\) or as the array
\(E=[e_1\ \cdots\ e_m]\). The expression \(v=Ea\) means
\[
  v=e_1a^1+\cdots+e_ma^m.
\]
The vector \(v\) belongs to \(V\); the column \(a\in\mathbb F^m\)
contains its coordinates in \(E\). Changing the basis changes the
coordinate column while leaving the vector itself unchanged.

The equation \(F=EM\) says that column \(j\) of \(M\) contains the
\(E\)-coordinates of \(f_j\). If \(v=Ea=Fb\), then \(a=Mb\), so
\(M\) converts \(F\)-coordinates to \(E\)-coordinates. See
\autoref{sec:linear-algebra-i:change-of-basis-for-vectors}.

For a linear map \(L:V\to W\), with domain basis \(E\) and codomain
basis \(F\), the notation
\[
  L(E)=[L(e_1)\ \cdots\ L(e_m)]=FM
\]
means that \(L\) is applied to each basis vector separately. Column
\(j\) of \(M\) contains the \(F\)-coordinates of \(L(e_j)\).
If \(\dim V=m\) and \(\dim W=n\), this matrix has size \(n\times m\).
The representing matrix depends on \(L\) and both ordered bases.
\autoref{tab:appendix:notation-maps} lists notation for maps.

\begin{table}[htbp]
  \centering
  \renewcommand{\arraystretch}{1.2}
\begin{tabular}{@{}p{0.23\linewidth}p{0.72\linewidth}@{}}
  \toprule
  Symbol & Meaning \\
  \midrule
  \(L:V\to W\) & \notationmeaning{A map with domain \(V\) and codomain \(W\).} \\
  \(v\mapsto L(v)\) & \notationmeaning{The rule assigning an output to an input.} \\
  \(\mathcal L(V,W)\) & \notationmeaning{The vector space of linear maps
    from \(V\) to \(W\).} \\
  \(I_V\) & \notationmeaning{The identity map on \(V\).} \\
  \(I_n\) & \notationmeaning{The \(n\)-by-\(n\) identity matrix.} \\
  \(L^{-1},M^{-1}\) & \notationmeaning{The inverse map and inverse matrix,
    respectively, when they exist.} \\
  \bottomrule
\end{tabular}
  \caption{Maps, identities, and inverses.}
  \label{tab:appendix:notation-maps}
\end{table}

Composition is read from right to left. In \(L_2\circ L_1\), apply
\(L_1\) first and \(L_2\) second; the representing matrix is
\(M_2M_1\) in compatible bases. A prime in \(E'\), \(F'\), or \(M'\)
labels a new basis or matrix, whereas \(f'\) denotes the derivative
of a function in the differentiation discussion.

\section{Kernel, image, and function spaces}
\label{sec:appendix:notation-kernel-and-image}

\begin{table}[htbp]
  \centering
  \renewcommand{\arraystretch}{1.2}
\begin{tabular}{@{}p{0.23\linewidth}p{0.72\linewidth}@{}}
  \toprule
  Symbol & Meaning \\
  \midrule
  \(\ker L\) & \notationmeaning{The subspace of inputs mapped to \(0_W\).} \\
  \(\operatorname{Null}(M)\) & \notationmeaning{The null space \(\ker M\)
    of a matrix.} \\
  \(\operatorname{im}L=L(V)\) & \notationmeaning{The subspace of outputs
    actually attained by \(L\).} \\
  \(\operatorname{nullity}L\) & \notationmeaning{The number \(\dim\ker L\).} \\
  \(\operatorname{rank}L\) & \notationmeaning{The number \(\dim\operatorname{im}L\).} \\
  \(v_0+\ker L\) & \notationmeaning{The set \(\{v_0+k\mid k\in\ker L\}\).} \\
  \bottomrule
\end{tabular}
  \caption{Kernel, image, rank, and nullity.}
  \label{tab:appendix:notation-kernel-and-image}
\end{table}

\autoref{tab:appendix:notation-kernel-and-image} lists kernel and image
(subspaces) and nullity and rank (their dimensions).
The codomain \(W\) specifies the target space, while the image records
which targets are reached. Writing \(\ker M\) or \(\operatorname{im}M\)
treats the matrix as the linear map \(x\mapsto Mx\).
See \autoref{sec:linear-algebra-i:kernel} through
\autoref{sec:linear-algebra-i:rank}.

Both \(C(\R)\) and \(C^0(\R)\) denote continuous real-valued functions
on \(\R\); \(C^1(\R)\) denotes continuously differentiable ones.
These superscripts specify regularity. The space \(\mathcal P\) contains
real polynomials of arbitrary finite degree; \(\mathcal P_d(\R)\) contains
those of degree at most \(d\), including the zero polynomial.
These vector spaces consist of functions; see
\autoref{sec:linear-algebra-i:differentiation}.

\section{Permutations and multilinear notation}
\label{sec:appendix:notation-permutations}

\begin{table}[htbp]
  \centering
  \renewcommand{\arraystretch}{1.2}
\begin{tabular}{@{}p{0.23\linewidth}p{0.72\linewidth}@{}}
  \toprule
  Symbol & Meaning \\
  \midrule
  \(S_n\) & \notationmeaning{Permutations of \(\{1,\ldots,n\}\),
    with composition.} \\
  \(T_n\) & \notationmeaning{All maps from \(\{1,\ldots,n\}\) to itself.} \\
  \(\iota\) & \notationmeaning{The identity permutation.} \\
  \(\tau_{ij}\) & \notationmeaning{The transposition exchanging \(i\) and \(j\).} \\
  \(\varepsilon(\sigma)\) & \notationmeaning{The sign of a permutation,
    equal to \(1\) or \(-1\).} \\
  \(\varepsilon(\phi)\) & \notationmeaning{The extended sign on \(T_n\),
    equal to \(0\) for a map that is not a permutation.} \\
  \bottomrule
\end{tabular}
  \caption{Permutations and their signs.}
  \label{tab:appendix:notation-permutations}
\end{table}

\autoref{tab:appendix:notation-permutations} uses \(T_n\) for a set of
maps, distinct from the transpose superscript \(T\).
As with maps, \(\sigma_2\circ\sigma_1\) applies \(\sigma_1\) first.
The indices in \(v_{\sigma(j)}\) and \(b^{\sigma(j)}_j\) select a vector
from the list and a row of column \(j\), respectively. See
\autoref{sec:linear-algebra-i:permutations} and
\autoref{sec:linear-algebra-i:determinant}.

\subsection{Symbols whose meaning depends on context}

The symbol \(A\) denotes a matrix in \(Ax=b\), oriented area in
\(A(v,w)\), or an antisymmetric multilinear function in
\(A(v_1,\ldots,v_n)\). Its inputs and domain identify the meaning.
The symbol \(\mathcal V(a,b,c)\) denotes oriented volume. Ordinary area
and volume are the absolute values of their oriented counterparts.

For differentiation, \(D(f)=f'\). For determinants, \(D\) is the
normalized antisymmetric multilinear function, with
\(\det M=D(C_1,\ldots,C_n)\) for the columns \(C_j\) of a square matrix.
These column labels are distinct from the field \(\C\) and the function
spaces \(C^0(\R)\) and \(C^1(\R)\).

The symbols \(P(v,w)\) and \(P(a,b,c)\) denote a parallelogram and a
parallelepiped, whereas \(\mathcal P\) is a polynomial space.
\endgroup
